\subsection{与对偶相容的拓扑}

本节中，E 与 F 表示由一个双线性型 B （II, 第 40 页）置于对偶的两个向量空间。我们回顾（II, 第 41 页）曾定义过两个线性映射

\[
d _ {\mathbf {B}}: \mathrm{F} \to \mathrm{E} ^ {*}, s _ {\mathbf {B}}: \mathrm{E} \to \mathrm{F} ^ {*}
\]

它们由关系式

\[
\mathbf {B} (x, y) = \langle x, d _ {\mathbf {B}} (y) \rangle = \langle y, s _ {\mathbf {B}} (x) \rangle\tag{1}
\]

刻画，其中 \(x\in \mathbf{E},y\in \mathbf{F}.\)

定义 1. --- E 上的局部凸拓扑 T 称为与 E 和 F 之间的对偶相容，如果 \(d_{B}\) 是从 F 到由赋以拓扑 T 的 E 所得局部凸空间的对偶上的一个双射。

如果存在这样一个拓扑 \(\mathcal{T}\) ，则映射 \(d_{\mathrm{B}}\) 是单射，也就是说，E 与 F 之间的对偶在 F 中是分离的（II, 第 41 页）。

命题 1. --- (i) E 中的闭凸子集对 E 上所有与 E 和 F 之间的对偶相容的局部凸拓扑来说是相同的。

\begin{enumerate}
\def\labelenumi{(\roman{enumi})}
\setcounter{enumi}{1}
\tightlist
\item
  \(\mathbf{E}\) 的有界子集对 \(\mathbf{E}\) 上所有与 \(\mathbf{E}\) 和 \(\mathbf{F}\) 之间的对偶相容的局部凸拓扑来说是相同的。
\end{enumerate}

设 T 是 E 上与 E 和 F 之间的对偶相容的一个拓扑，因而比 \(\sigma(\mathrm{E},\mathrm{F})\) 细。若 E 的一个凸子集对 T 是闭的，则它是闭的实半空间之交（II, 第 38 页, 推论 1），因而它对 \(\sigma(\mathrm{E},\mathrm{F})\) 是闭的。这就证明了 (i)。断言 (ii) 已在 III, 第 27 页, 推论 3 中证明。

用 \(F_{\sigma}\) 表示赋以弱拓扑 \(\sigma(F, E)\) 的向量空间 F。那么线性映射 \(s_{B}\) 把 E 映到对偶 \((F_{\sigma})'\) （\(F_{\sigma}\) 的）上（II, 第 43 页, 命题 3）。设 \(S\) 是 \(F_{\sigma}\) 的有界子集的一个族。滥用术语，我们把在 \(s_{B}\) 之下的 \(S\) -拓扑（\((F_{\sigma})'\) 上的）的原像称为 E 上的 \(S\) -拓扑。它由半范数族

\[
p _ {\mathbf {A}} (x) = \sup _ {y \in \mathbf {A}} | \mathrm{B} (x, y) |,\tag{2}
\]

所定义，其中 A 取遍 S。特别地，当 S 是 F 的有限子集的族时，S-拓扑正是弱拓扑 \(\sigma(\mathrm{E}, \mathrm{F})\) 。

定义 2. --- 设 E 与 F 是两个处于对偶的空间。E 上的 Mackey 拓扑，记作 \(\tau(\mathrm{E},\mathrm{F})\) ，定义为 E 上的 \(\mathfrak{S}\) -拓扑，其中 \(\mathfrak{S}\) 是 F 的这样一些子集全体所成的族：它们在 \(E^{*}\) 中的像（在 \(d_{\mathrm{B}}\) 之下）对 \(\sigma(\mathrm{E}^{*},\mathrm{E})\) 是凸的、均衡的且紧的。

当 E 与 F 之间的对偶在 F 中是分离的时，\(d_{B}\) 是单射，且 F 上的拓扑 \(\sigma(F, E)\) 是在 \(d_{B}\) 之下的拓扑 \(\sigma(E^{*}, E)\) （\(E^{*}\) 上的）的原像。在这种情形，S 由 F 的对 \(\sigma(F, E)\) 为凸的、均衡的且紧的那些子集组成。

一般地，若 \(F_{1}=d_{\mathrm{B}}(\mathrm{F})\subset\mathrm{E}^{*}\) ，并且用 \((x,y_{1})\mapsto\mathrm{B}_{1}(x,y_{1})\) 表示典范双线性型 \((x,x^{*})\mapsto\langle x,x^{*}\rangle\) 在 \(E\times F_{1}\) 上的限制，则 E 与 \(F_{1}\) 由 \(B_{1}\) 置于对偶，且这个对偶在 \(F_{1}\) 中是分离的，因为按定义有 \(\mathrm{B}(x,y)=B_{1}(x,d_{\mathrm{B}}(y))\) ，定义 2 表明 \(\tau(\mathrm{E},\mathrm{F})=\tau(\mathrm{E},\mathrm{F}_{1})\) 。

注记 1. --- 设 A 是 Hausdorff 局部凸空间 G 的一个紧凸子集，\(\tilde{A}\) 是 A 的闭均衡凸包。当域 K 为 R 时，集 \(\tilde{A}\) 是 \(A \cup (-A)\) 的闭凸包；当 K 为 C 时，集 \(\tilde{A}\) 含于 \(2A \cup (-2A) \cup (2iA) \cup (-2iA)\) 的闭凸包。因此（II, 第 14 页, 命题 15），\(\tilde{A}\) 是紧的。

特别地，由此推出，当 E 与 F 之间的对偶在 F 中是分离的时，Mackey 拓扑 \(\tau(\mathrm{E},\mathrm{F})\) 也是 \(S'\) -拓扑，其中 \(S'\) 是 F 的对 \(\sigma(\mathrm{F},\mathrm{E})\) 为紧的所有凸子集所成的集。以类似的方式我们在 F 上定义 Mackey 拓扑 \(\tau(\mathrm{F},\mathrm{E})\) 。

定理 1（Mackey）. --- 设 E 与 F 是两个处于对偶的空间；假设这个对偶在 F 中是分离的。要使 E 上的局部凸拓扑 T 与 E 和 F 之间的对偶相容，必须且只需 T 比拓扑 \(\sigma(\mathrm{E},\mathrm{F})\) 细且比 Mackey 拓扑 \(\tau(\mathrm{E},\mathrm{F})\) 粗。

把 F 与它在 \(E^{*}\) 中的像（在 \(d_{B}\) 之下）等同。用 \(S_{0}\) 表示 F 的对 \(\sigma(F, E)\) 为凸的、均衡的且紧的所有子集所成的集。按定义，\(\tau(E, F)\) 是 E 上的 \(S_{0}\) -拓扑，因而比 \(\sigma(E, F)\) 细。

引理 1. --- E* 的子空间 F 由 E 上对 \(\tau(\mathrm{E},\mathrm{F})\) 连续的所有线性型组成。

F 的每个元素都是对 \(\sigma(\mathrm{E},\mathrm{F})\) 连续的映射，因而对 \(\tau(\mathrm{E},\mathrm{F})\) 连续。

反之，设 \(f \in \mathrm{E}^*\) 对 \(\tau(\mathrm{E}, \mathrm{F})\) 连续。存在 0 的一个邻域 \(\mathbf{U}\) （\(\mathrm{E}\) 中，对 \(\tau(\mathrm{E}, \mathrm{F})\) 而言），使得 \(|f| \leqslant 1\) 在 \(\mathbf{U}\) 上成立；可以假设存在一个集 \(\mathbf{A} \in \mathfrak{S}_0\) 使得 \(\mathbf{U} = \mathbf{A}^\circ\) 。换言之，\(f\) 属于双极集 \(\mathbf{A}^{\circ\circ}\) （\(\mathbf{A}\) 关于 \(\mathbf{E}^*\) 与 \(\mathbf{E}\) 之间的对偶而言）。但拓扑 \(\sigma(\mathrm{F}, \mathrm{E})\) （\(\mathbf{F}\) 上的）由 \(\sigma(\mathrm{E}^*, \mathrm{E})\) 诱导；于是 \(\mathbf{A}\) 对 \(\sigma(\mathrm{E}^*, \mathrm{E})\) 是凸的、均衡的且紧的，且双极定理（II, 第 44 页, 定理 1）蕴含等式 \(\mathbf{A} = \mathbf{A}^{\circ\circ}\) 。因此有 \(f \in \mathrm{F}\) ，引理由此得证。

引理 2. --- 设 T 是 E 上的一个局部凸拓扑，使得 E 上对 T 连续的每个线性型都属于 F。那么 T 比 \(\tau(\mathrm{E}, \mathrm{F})\) 粗。

设 \(\mathfrak{U}\) 是对 \(\mathcal{T}\) 而言 0 的凸均衡邻域所成的集。设 \(\mathfrak{S}\) 是 \(\mathfrak{U}\) 中诸元素在 F 中的极集所成的集。由 III, 第 17 页, 推论 2，有 \(\mathfrak{S} \subset \mathfrak{S}_0\) ；又由 III, 第 19 页, 命题 7 的推论 1，\(\mathcal{T}\) 与 \(\mathfrak{S}'\) -拓扑恒同，其中 \(\mathfrak{S}'\) 是 \(\mathfrak{U}\) 中诸集在 E 的对偶 E' 中的极集所成的集。但由假设 E' \(\subset\) F，因而 \(\mathfrak{S}'\) 的每个集都含于 \(\mathfrak{S}\) 的某个集中；引理由此得证。

设 T 是 E 上与 E 和 F 之间的对偶相容的一个拓扑。由引理 2，T 比 \(\tau(\mathrm{E},\mathrm{F})\) 粗，且显然 T 比 \(\sigma(\mathrm{E},\mathrm{F})\) 细。反之，对拓扑 \(\tau(\mathrm{E},\mathrm{F})\) （引理 1）以及对拓扑 \(\sigma(\mathrm{E},\mathrm{F})\) （II, 第 43 页, 命题 3），F 都是 E 的对偶，因而对介于 \(\tau(\mathrm{E},\mathrm{F})\) 与 \(\sigma(\mathrm{E},\mathrm{F})\) 之间的每个拓扑，F 也是 E 的对偶。

推论. --- 设 p 是 E 上的一个半范数。下列条件等价：

\begin{enumerate}
\def\labelenumi{(\roman{enumi})}
\item
  \(p\) 对拓扑 \(\tau(\mathrm{E},\mathrm{F})\) 连续；
\item
  每个线性型 \(f\) （\(\mathbf{E}\) 上的）若满足 \(|f| \leqslant p\) ，都来自 \(\mathbf{F}\) 的一个元素。
\item[]
  \((i) \Rightarrow (ii):\) 若 \(p\) 对 \(\tau(\mathrm{E},\mathrm{F})\) 连续，则每个线性型 \(f\) （\(\mathrm{E}\) 上的）若满足 \(|f|\leqslant p\) ，就对 \(\tau(\mathrm{E},\mathrm{F})\) 连续，因而由引理 1 来自 \(\mathrm{F}\) 的一个元素。
\item[]
  \((ii) \Rightarrow (i):\) 设 T 是 E 上由半范数 p.~定义的拓扑。若条件 (ii) 满足，则 E 上对 T 连续的线性型都属于 F。由引理 2，T 比 \(\tau(\mathrm{E}, \mathrm{F})\) 粗，因而 p 对 \(\tau(\mathrm{E}, \mathrm{F})\) 连续。
\end{enumerate}

注记 2. --- * 设 K 是 F 的一个对弱拓扑 \(\sigma(F, E)\) 紧的凸子集，\(\mu\) 是 K 上的一个正测度。令

\[
p (x) = \int_ {\mathbf {K}} | \mathrm{B} (x, y) | d \mu (y)
\]

（对一切 \(x \in E\) ）。立即可见 p 是一个半范数。此外，对每个 \(x \in E\) ，关系式 \(\langle |B(x, y)| \leqslant 1 \text{ 对一切 } y \in K \rangle\) 蕴含 \(p(x) \leqslant \mu(K)\) 。这就证明 E 上的半范数 p 对 Mackey 拓扑 \(\tau(E, F)\) 连续。*

例. --- 设 G 是一个局部凸空间，\(G'\) 是它的对偶。在 \(G'\) 上，弱拓扑 \(\sigma(G', G)\) 与紧凸收敛拓扑（III, 第 14 页）都与 \(G'\) 和 G 之间的对偶相容。一般说来，\(G'\) 上的强拓扑与紧收敛拓扑都不与 \(G'\) 和 G 之间的对偶相容。然而回想，当 G 是 Hausdorff 的且拟完备的时候，\(G'\) 上的紧收敛拓扑与紧凸收敛拓扑一致（III, 第 8 页），因而与 \(G'\) 和 G 之间的对偶相容。

定义 3. --- 设 E 与 F 是两个处于对偶的向量空间，T 是 F 的对 \(\sigma(\mathrm{F}, \mathrm{E})\) 有界的子集所成的族。那么 F 上的 T-拓扑记作 \(\beta(\mathrm{E}, \mathrm{F})\) 。

类似地，我们在 F 上定义拓扑 \(\beta(\mathrm{F},\mathrm{E})\) 。容易看出，拓扑 \(\beta(\mathrm{E},\mathrm{F})\) 与 \(\beta(\mathrm{E},\mathrm{F}/\mathrm{E}^{\circ})\) 恒同，因而可以化归为 E 与 F 之间的对偶在 F 中是分离的情形。

注记. --- 3) 用 \(E_{\sigma}\) 表示赋以拓扑 \(\sigma(\mathrm{E}, \mathrm{F})\) 的空间 E。\(E_{\sigma}\) 中的桶（III, 第 24 页）是 E 的对 \(\sigma(\mathrm{E}, \mathrm{F})\) 为凸的、均衡的、闭的且吸收的子集。它们正是 F 的对 \(\sigma(\mathrm{F}, \mathrm{E})\) 为凸的、均衡的且有界的子集的极集。因此，\(E_{\sigma}\) 中所有桶的族是 E 中拓扑 \(\beta(\mathrm{E}, \mathrm{F})\) 的 0 的一个基本邻域系。换言之，E 上的一个半范数对 \(\beta(\mathrm{E}, \mathrm{F})\) 连续，当且仅当它对 \(\sigma(\mathrm{E}, \mathrm{F})\) 下半连续（cf.~III, 第 24 页, 命题 1）。

\begin{enumerate}
\def\labelenumi{\arabic{enumi})}
\setcounter{enumi}{3}
\item
  设 \(\mathcal{T}\) 是 E 上与 E 和 F 之间的对偶相容的一个拓扑。由 IV, 第 1 页, 命题 1, (ii)，F 上的拓扑 \(\beta(\mathrm{F},\mathrm{E})\) 正是 F 上的强拓扑（当把 F 与 E 的对偶（赋以拓扑 \(\mathcal{T}\) ）认同时）。
\item
  拓扑 \(\beta(\mathrm{E},\mathrm{F})\) （\(\mathrm{E}\) 上的）比 \(\tau(\mathrm{E},\mathrm{F})\) 细。一般说来，它与 \(\mathrm{E}\) 和 \(\mathrm{F}\) 之间的对偶不相容（但 cf.~§ 2）。特别地，\(\mathrm{E}\) 的一个对 \(\sigma(\mathrm{E},\mathrm{F})\) 有界的子集不一定对 \(\beta(\mathrm{E},\mathrm{F})\) 有界。
\end{enumerate}
% semantic-fragments: legacy-input-seam
\relax{}
